% Part: turing-machines
% Chapter: undecidability
% Section: trakhtenbrot

\documentclass[../../../include/open-logic-section]{subfiles}

\begin{document}

\olfileid{tur}{und}{tra}
\olsection{د تراختن‌بروت قضيه}

\begin{explain}
  په~\olref[rep]{sec} کښې مو د يوه ټيورينګ ماشين~$M$ او ننوت
  رشته~$w$ دپاره !!{sentence}s~$!T(M,w)$ او~$!E(M,w)$ تعريف کړل۔
  بيا مو په~\olref[ver]{lem:valid-if-halt} او
  \olref[ver]{lem:halt-if-valid} کښې وښودل چې~$!T(M,w) \lif
  !E(M,w)$ هله او يوازې هله معتبر دے چې~$M$ پر ننوت~$w$ له پيلېدو
  وروسته بالاخره ودرېږي۔ څرنګه چې د درېدنې مسئله پرېکړه کېدونکې نۀ
  ده، له دې څخه پايله اخلو چې د لومړۍ درجې منطق د !!{sentence}s
  اعتبار او د صدق وړتيا پرېکړه کېدونکي نۀ دي
  (\cref{tur:und:uns:thm:decision-prob,tur:und:uns:cor:undecidable-sat})۔

  خو د جملو اعتبار او د صدق وړتيا د خپلو‌سرو
  !!{structure}s دپاره تعريف شوي، که هغه متناهي وي او که نامتناهي۔
  ښايي تاسو ګومان وکړئ چې دا پرېکړه اسانه ده چې آيا يو
  !!a{sentence} په يوه متناهي !!{structure} کښې د صدق وړ دے (يا په
  ټولو متناهي !!{structure}s کښې معتبر دے)۔ د پرېکړې مسئلې د ناحلي
  ثبوت داسې برابرولے شو چې وښيي خبره داسې نۀ ده۔

  لومړے، که د~\olref[ver]{lem:halt-if-valid} ثبوت ته بېرته ورشئ،
  وبه وينئ چې هلته مو د~$!T(M,w)$ داسې مدل~$\Struct M$ جوړ کړ چې
  بالکل دا بيانوي چې ماشين~$M$ پر ننوت~$w$ له پيلېدو وروسته څه
  کوي۔ د هغه مدل ډومېن~$\Nat$ و، يعنې نامتناهي و۔ خو که~$M$ په
  رښتيا پر ننوت~$w$ ودرېږي، په همدې ډول يو متناهي مدل~$\Struct M'$
  جوړولے شو۔ فرض کړئ~$M$ پر ننوت~$w$ له پيلېدو وروسته له~$k$
  ګامونو وروسته درېږي۔ د ډومېن~$\Domain{M'}$ دپاره سټ~$\{0,
  \dots, n\}$ واخلئ، چې~$n$ د~$k$ او د~$w$ د اوږدوالي تر منځ لوے
  عدد دے، او داسې وټاکئ:
  \[
    \Assign{\prime}{M'}(x) =
    \begin{cases}
      x + 1 &\text{که $x < n$ وي}\\
      n &\text{که نۀ وي،}
    \end{cases}
  \]
  او~$\tuple{x,y} \in \Assign{<}{M'}$ هله او يوازې هله چې~$x < y$
  يا~$x = y = n$ وي۔ له دې پرته~$\Struct{M'}$ بالکل د~$\Struct{M}$
  په شان تعريفېږي۔ د~$\Struct{M'}$ د تعريف له مخې، لکه د
  \olref[ver]{lem:halt-if-valid} په ثبوت کښې،
  $\Sat{M'}{!T(M,w)}$۔ او څرنګه چې مو فرض کړې ده چې~$M$ پر
  ننوت~$w$ درېږي، $\Sat{M'}{!E(M,w)}$۔ نو~$\Struct{M'}$ د~$!T(M,w)
  \land !E(M,w)$ يو متناهي مدل دے (پام وکړئ چې~$\lif$ مو
  په~$\land$ بدل کړے دے)۔

  موږ د ثبوت نيمې لارې ته رسېدلي يو: ښودلې مو ده چې که~$M$ پر
  ننوت~$w$ ودرېږي، نو~$!T(M,w) \land !E(M,w)$ يو متناهي مدل لري۔
  له بده مرغه د دې نقيض عکس سم نۀ دے؛ يعنې داسې ټيورينګ ماشينونه
  شته چې پر کوم ننوت~$w$ نۀ درېږي، خو~$!T(M,w) \land !E(M,w)$ بيا
  هم يو متناهي مدل لري۔ د بېلګې په توګه هغه ماشين~$M$ په پام کښې
  ونيسئ چې يوازې يو حالت~$q_0$ او دا لارښوونه لري:
  $\delta(q_0,\TMblank) = \tuple{q_0,\TMblank,\TMstay}$۔ دا ماشين
  پر تش ننوت~$w = \emptyseq$ له پيلېدو وروسته هېڅکله نۀ درېږي: په
  يوه نامتناهي کړۍ کښې دے، خو پټه نۀ بدلوي او سر نۀ خوځوي۔ ټول
  تشکيلونه يو شان دي (هماغه حالت، هماغه د سر موقعيت، هماغه د پټې
  منځپانګه)۔ داسې متناهي !!{structure}~$\Struct{M''}$ تعريفولے شو
  چې~$!T(M,\emptyseq) \land !E(M,\emptyseq)$ پوره کړي (تمرين)۔ خو
  $!T(M,w)$ په مناسبه توګه بدلولے شو، څو داسې !!{structure}s ترې
  بهر کړو۔

  \begin{prob}
    فرض کړئ~$M$ داسې ټيورينګ ماشين دے چې يوازې حالت~$q_0$ او يوازې
    لارښوونه~$\delta(q_0,\TMblank) =
    \tuple{q_0,\TMblank,\TMstay}$ لري۔ داسې وټاکئ:
    $\Domain{M''} = \{0, 1, 2\}$، $\Assign{\prime}{M''}(0) =
    \Assign{\prime}{M''}(1) = 1$ او~$\Assign{\prime}{M''}(2) = 2$،
    او~$\Assign{<}{M''} = \{\tuple{0,1}, \tuple{1,1},
    \tuple{2,2}\}$۔ $\Assign{\Obj Q_{q_0}}{M''}$، $\Assign{\Obj
    S_{\TMblank}}{M''}$، او~$\Assign{\Obj S_{\TMendtape}}{M''}$
    داسې تعريف کړئ چې~$!T(M,\emptyseq)$ او~$!E(M,\emptyseq)$ رښتوني
    شي، او څرګنده کړئ چې ولې رښتوني دي۔ اشاره: پام وکړئ چې
    $\delta(q_0, \TMendtape)$ ناتعريفه ده۔ ډاډ ترلاسه کړئ چې
    \begin{align*}
    & \Obj Q_{q_0}(\num{1}, \num{n}) \land \Obj S_{\TMendtape}(\num{0}, \num{n})
      \land
      \lforall[x][(\num{0} < x \lif \Obj S_\TMblank(x, \num{n}))] \text{\quad د هر $n \in \Nat$ دپاره}\\
    & \lexists[y][(\Obj Q_{q_0}(\num{0}, y) \land \Obj S_{\TMendtape}(\num{0}, y))]
    \end{align*}
    دواړه په~$\Struct{M''}$ کښې رښتوني وي۔
    \textbf{د ژباړې د نسخې سمون (OLCMP-063):} سرچينې د يوازېني
    حالت پر ځاے په لارښوونه کښې يو ناتعريفه حالت راوړے و؛ دلته هماغه
    ټاکل شوے يوازېنی حالت کارول شوے دے۔
    \textbf{د ژباړې د نسخې سمون (OLCMP-064):} سرچينې د يو د تالي
    ارزښت په تېروتنه د پخواني جوړښت په نوم پورې تړلے و؛ دلته د تمرين
    روان جوړښت په يو شان ډول کارول شوے دے۔
  \end{prob}
\end{explain}

هغه !!{sentence}s په پام کښې ونيسئ چې د ټيورينګ ماشين~$M$ عمل پر
ننوت~$w = \sigma_{i_1}\dots\sigma_{i_k}$ بيانوي:
\begin{enumerate}
\item هغه اکسيومونه چې اعداد او~$<$ بيانوي (بالکل لکه په
\olref[rep]{sec} کښې د~$!T(M,w)$ په تعريف کښې)۔
\item هغه اکسيومونه چې د ننوت تشکيل بيانوي: بالکل لکه د~$!T(M,w)$
په تعريف کښې۔
\item هغه اکسيومونه چې له يوه تشکيل څخه بل ته د حالت بدلون بيانوي:

د لاندې برخې دپاره، $!A(x, y)$ دې د پخوا په شان وي، او داسې دې وي:
\[
  !B(y) \ident \lforall[x][(x < y \lif \eq/[x][y])].
\]
\begin{enumerate}
\item \ollabel{rep-right} د هرې لارښوونې~$\delta(q_i, \sigma) =
  \tuple{q_j, \sigma', \TMright}$ دپاره دا !!{sentence}:
\begin{align*}
& \lforall[x][\lforall[y][(
   (\Obj Q_{q_i}(x, y) \land \Obj S_{\sigma}(x, y)) \lif {}]] \\
&\qquad   (\Obj Q_{q_j}(x', y') \land \Obj S_{\sigma'}(x, y') \land
!A(x, y) \land !B(y')))
\end{align*}
\item \ollabel{rep-left} د هرې لارښوونې~$\delta(q_i, \sigma) =
  \tuple{q_j, \sigma', \TMleft}$ دپاره دا !!{sentence}:
\begin{align*}
& \lforall[x][\lforall[y][
    ((\Obj Q_{q_i}(x', y) \land \Obj S_{\sigma}(x', y)) \lif {}]]\\
& \qquad   (\Obj Q_{q_j}(x, y') \land \Obj S_{\sigma'}(x', y') \land
!A(x', y) \land !B(y'))) \land {}\\
& \lforall[y][((\Obj Q_{q_i}(\num{0}, y) \land \Obj S_{\sigma}(\num{0},
    y)) \lif {}]\\
& \qquad (\Obj Q_{q_j}(\num{0}, y') \land \Obj S_{\sigma'}(\num{0},
  y') \land !A(\num{0}, y) \land !B(y')))
\end{align*}
\textbf{د ژباړې د نسخې سمون (OLCMP-065):} د سرچينې د کيڼ حرکت
عادي څانګې بيا د نابدلې پټې په شرط کښې د ليکل کېدونکې مربع پر ځاے
منزل مستثنا کاوه؛ دلته د ليکل کېدونکې مربع سم آرګومان کارول شوے دے۔
\textbf{د ژباړې د نسخې سمون (OLCMP-066):} د سرچينې همدې عادي څانګې
هغه د بېلوالي شرط پرېښے و چې د دې درې واړو حالت-بدلونونو موخه ده؛
دلته ور زيات شوے، څو د بل تشکيل عدد له ټولو مخکښېنيو بېل وي۔
\item \ollabel{rep-stay} د هرې لارښوونې~$\delta(q_i, \sigma) =
  \tuple{q_j, \sigma', \TMstay}$ دپاره دا !!{sentence}:
\begin{align*}
& \lforall[x][\lforall[y][(
   (\Obj Q_{q_i}(x, y) \land \Obj S_{\sigma}(x, y)) \lif {}]] \\
&\qquad (\Obj Q_{q_j}(x, y') \land \Obj S_{\sigma'}(x, y') \land
!A(x, y) \land !B(y')))
\end{align*}
\end{enumerate}
لکه چې وينئ، هغه !!{sentence}s چې د~$M$ حالت-بدلونونه بيانوي د
$!T(M,w)$ له اړوندو !!{sentence} سره يو شان دي، يوازې په پای کښې
$!B(y')$ ور زياتوو۔ $!B(y')$ دا تضمينوي چې د ``بل'' تشکيل عدد~$y'$
له ټولو مخکښېنيو عددونو~$\Obj 0$، $\Obj 0'$، \dots، څخه بېل وي۔
% \item Sentences that express consistency conditions:
% \begin{enumerate}
%   \item\ollabel{rep-con-symb}%
%   !!^a{sentence} that says that no square can have more than one
%   symbol on it at any given time:
%   \[\lforall[x][\lforall[y][(\Obj S_{\sigma}(x, y) \lif \lnot \Obj
%   S_{\sigma'}(x, y))]],\]
%   for every pair $\sigma \neq \sigma'$ of tape symbols of~$T$.
%   \item\ollabel{rep-con-state}%
%   !!^a{sentence} that says that $M$ cannot be in more than one state
%   at any given time:
%   \[\lforall[x][\lforall[y][(\Obj Q_{q}(x, y) \lif
%   \lforall[z][\lnot \Obj
%   Q_{q'}(z, y))]],\]
%   for every pair $q \neq q'$ of states of~$T$.
%   \item\ollabel{rep-con-tape}%
%   !!^a{sentence} that says that $M$ cannot be on more than one tape
%   square at any given time:
%   \[\lforall[x][\lforall[y][(\Obj Q_{q}(x, y) \lif
%   \lforall[z][\eq/[z][y]\lif \lnot \Obj
%   Q_{q'}(z, y))]],\]
%   for every pair $q$, $q'$ of states of~$T$.
% \end{enumerate}
\end{enumerate}
$!T'(M, w)$ د ټيورينګ ماشين~$M$ او ننوت~$w$ دپاره د پورتنيو ټولو
!!{sentence}s عطف وي۔

\begin{lem}\ollabel{lem:halts-sat}
  که~$M$ پر ننوت~$w$ له پيلېدو وروسته ودرېږي، نو~$!T'(M,w) \land
  !E(M,w)$ يو متناهي مدل لري۔
\end{lem}

\begin{proof}
  $\Struct{M'}$ د~\olref[ver]{lem:halt-if-valid} په ثبوت کښې د
  جوړښت په شان واخلئ، خو دا توپيرونه ولري:
  \begin{align*}
    \Domain{M'} & = \{0, \dots, n\},\\
    \Assign{\prime}{M'}(x) & =
    \begin{cases}
      x + 1 &\text{که $x < n$ وي}\\
      n &\text{که نۀ وي،}
    \end{cases} \\
    \tuple{x,y} \in \Assign{<}{M'} &\text{هله او يوازې هله چې $x < y$ يا $x = y = n$ وي،}
  \end{align*}
  چې~$n = \max(k,\len{w})$ او~$k$ تر ټولو کوچنی داسې عدد دے چې~$M$
  پر ننوت~$w$ له پيلېدو وروسته له~$k$ ګامونو وروسته درېدلے وي۔ د دې
  خبرې تاييد د تمرين په توګه پرېږدو چې~$\Sat{M'}{!T'(M,w) \land
  !E(M,w)}$۔
  \textbf{د ژباړې د نسخې سمون (OLCMP-067):} سرچينې په وروستي
  پوره-کېدو حکم کښې د درېدنې ټاکلې څيز-ژبې نښه پرېښې وه؛ دلته بشپړه
  جمله راوړل شوې ده۔
\end{proof}

\begin{prob}
  د~\olref[tur][und][tra]{lem:halts-sat} ثبوت په دې ښودلو بشپړ کړئ
  چې~$\Sat{M'}{!T'(M,w) \land !E(M,w)}$۔
  \textbf{د ژباړې د نسخې سمون (OLCMP-068):} سرچينې په دې تمرين کښې
  د بدلې شوې تمثيلي جملې پر ځاے پخوانۍ جمله راوړې وه او د درېدنې د
  جملې ټاکلې نښه ئې پرېښې وه؛ دلته دواړه د لمې له بيان سره سم شوي دي۔
\end{prob}

\begin{lem}\ollabel{lem:sat-halts}
  که~$!T'(M,w) \land !E(M,w)$ يو متناهي مدل ولري، نو~$M$ پر
  ننوت~$w$ له پيلېدو وروسته درېږي۔
\end{lem}

\begin{proof}
  نقيض عکس ښيو۔ فرض کړئ~$M$ پر ننوت~$w$ له پيلېدو وروسته نۀ
  درېږي۔ که~$!T'(M,w) \land !E(M,w)$ هېڅ مدل ونۀ لري، خبره بشپړه
  ده۔ نو فرض کړئ~$\Struct{M}$ د~$!T'(M,w) \land !E(M, w)$ يو مدل
  دے۔ بايد وښيو چې دا متناهي کېدے نۀ شي۔
  \textbf{د ژباړې د نسخې سمون (OLCMP-069):} سرچينې دلته د روانې
  بدلې شوې جملې پر ځاے پخوانۍ تمثيلي جمله د مدل فرض ګرځولې وه؛ دلته
  فرض د لمې او د پاتې ثبوت له بدلې شوې جملې سره يو شان شوے دے۔

  بالکل د~\olref[ver]{lem:config} په شان ثابتولے شو چې که~$M$ پر
  ننوت~$w$ له پيلېدو وروسته تر~$n$ ګامونو پورې نۀ وي درېدلے، نو
  $!T'(M,w) \Entails !C(M, w, n) \land !B(\num{n})$۔ څرنګه چې~$M$
  پر ننوت~$w$ له پيلېدو وروسته نۀ درېږي، د هر~$n \in \Nat$ دپاره
  $!T'(M,w) \Entails !C(M, w, n) \land !B(\num{n})$۔ پام وکړئ چې
  د~\olref[rep]{prop:mlessk} له مخې، د ټولو~$k < n$ دپاره~$!T'(M,w)
  \Entails \num{k} < \num{n}$۔ همدارنګه~$!B(\num{n}) \Entails
  \num{k} < \num{n} \lif \eq/[\num{k}][\num{n}]$۔ نو د ټولو~$k <
  n$ دپاره~$\Sat{M}{\eq/[\num{k}][\num{n}]}$؛ يعنې نامتناهي ډېر
  ترمونه~$\num{k}$ بايد په~$\Struct{M}$ کښې ټول بېلابېل ارزښتونه
  ولري۔ خو دا غوښتنه کوي چې~$\Domain{M}$ نامتناهي وي؛ نو
  $\Struct{M}$ د~$!T'(M,w) \land !E(M, w)$ متناهي مدل نۀ شي
  کېدے۔
\end{proof}

\begin{prob}
  د~\olref[tur][und][tra]{lem:sat-halts} ثبوت په دې ښودلو بشپړ کړئ
  چې که~$M$ پر ننوت~$w$ له پيلېدو وروسته تر~$n$ ګامونو پورې نۀ وي
  درېدلے، نو~$!T'(M,w) \Entails !B(\num{n})$۔
\end{prob}

\begin{thm}[د تراختن‌بروت قضيه]
  \ollabel{thm:trakhtenbrodt}
  دا پرېکړه کېدونکې نۀ ده چې آيا د لومړۍ درجې منطق کوم خپل‌سری
  !!{sentence} يو متناهي مدل لري (يعنې په يوه متناهي جوړښت کښې د
  صدق وړ دے)۔
\end{thm}

\begin{proof}
  فرض کړئ داسې ټيورينګ ماشين~$F$ وے چې په يوه متناهي جوړښت کښې د
  صدق وړتيا پرېکړه کوي۔ بيا به مو د هر ټيورينګ ماشين~$M$ او
  ننوت~$w$ دپاره جمله~$!T'(M,w) \land !E(M,w)$ محاسبه کړې وے، او
  د~$F$ په وسيله به مو پرېکړه کړې وے چې آيا دا يو متناهي مدل لري۔
  د~\cref{tur:und:tra:lem:halts-sat,tur:und:tra:lem:sat-halts} له
  مخې، دا هله او يوازې هله متناهي مدل لري چې~$M$ پر ننوت~$w$ له
  پيلېدو وروسته ودرېږي۔ نو د~$F$ په وسيله به مو د درېدنې مسئله حل
  کړې وے، حال دا چې پوهېږو دا ناحل ده۔
\end{proof}

\begin{cor}\ollabel{cor:fproof-incomp}%
  داسې هېڅ !!{derivation} نظام نشي موجودېدے چې د \emph{متناهي}
  اعتبار دپاره سم او بشپړ وي؛ يعنې داسې !!a{derivation} نظام چې
  $\Proves !B$ هله او يوازې هله ولري چې د هرې متناهي
  !!{structure}~$\Struct{M}$ دپاره~$\Sat{M}{!B}$ وي۔
\end{cor}

\begin{proof}
  تمرين۔
\end{proof}

\begin{prob}
  \olref[tur][und][tra]{cor:fproof-incomp} ثابت کړئ۔ پام وکړئ چې~$!B$
  په هره متناهي !!{structure} کښې هله او يوازې هله پوره کېږي چې
  $\lnot !B$ په يوه متناهي جوړښت کښې د صدق وړ نۀ وي۔ څرګنده کړئ چې
  ولې په يوه متناهي جوړښت کښې د صدق وړتيا د
  \olref[tur][und][uns]{thm:valid-ce} په معنا نيمه د پرېکړې وړ ده۔
  له دې څخه دا استدلال وکړئ چې که د متناهي اعتبار دپاره کوم
  !!a{derivation} نظام وے، نو په يوه متناهي جوړښت کښې د صدق وړتيا
  به پرېکړه کېدونکې وه۔
\end{prob}


\end{document}
